\section{Introduction}

\citet{chollet2019measure} argued that intelligence should be measured by how efficiently a system turns priors and experience into skill, on tasks that lie largely outside its experience and have a verifiable answer. ARC is the concrete instance of that argument, and on its current edition the gap between what a person solves in a minute and what a model solves from the same demonstrations is still significant \citep{chollet2026arcprize2025}. For agent systems that continually operate for days and weeks, a more pressing question is whether a skill acquired on one job is still there the next time a similar scenario is met, which we call memorization, and whether that experience makes the system better at similar jobs it has never met, which we call generalization. By design, ARC measures acquisition once, within a single episode, and does not ask either over longer horizons.

The benchmarks that do ask it have three limitations (Section~\ref{sec:related}). Firstly, each is built around one \emph{applied} domain, such as coding or office work \citep{wu2024streambench,zheng2025lifelongagentbench,joshi2025swebenchcl}, in which the models under test already know a great deal from the pretraining data, and gains may come partly from drawing out that knowledge. Secondly, most serve a fixed set of instances, each seen once, so the same conceptual task never returns with a fresh input. This means a memorized answer cannot be told apart from a more abstract learned skill. Finally, while some permit both \citep{wu2024streambench,zheng2025lifelongagentbench}, most assume a single family of learner, gradient updates for fine-tuning suites such as MLLM-CTBench \citep{guo2025mllmctbench} and prompting and memory for agent benchmarks such as AgentCL and PATH-Bench \citep{shu2026agentcl,yang2026pathbench}.

We propose \bench{}, a lifelong stream built from ARC, where tasks need no knowledge of any domain beyond core-knowledge priors, such as objectness and counting \citep{chollet2019measure}, and the rule of each task has to be induced from a few demonstrations. The models we test solve \emph{some} of these tasks at first sight, but far from all (Appendix~\ref{app:tracks}). Continual-ARC has a track for each of ARC's three editions, which rise in difficulty \citep{chollet2025arcagi2,arcprize2026arcagi3}, to the point where cheaper models and frontier models are best measured on different tracks. For our ARC-AGI-1 track, the Re-ARC generators \citep{hodel2024rearc} produce a fresh, verified instance of any ARC-AGI-1 training task on demand. A task then returns after a dormant gap as often as the stream asks, each time with a new input that no cached answer fits. New tasks also keep arriving throughout the stream. On each instance, the system may submit an answer and receive a bare verdict (paying for incorrect attempts), pay for a demonstration, or take actions inside the harness if one is provided. We score the instance by the share of this feedback budget it used, all of it if the instance fails, as \citet{chollet2019measure} proposed for an interactive version of ARC. After the stream, we freeze the system and test it on fresh instances of the stream's tasks, which measures memorization, and on 100 tasks from ARC-AGI-1's evaluation set that the stream never shows, which measures generalization. Because a system may carry anything between instances, prompting, memory and weight updates all face the same stream and the same tests.

Most of our experiments focus on one agent harness, prime-agent \citep{karten2026primeagent}, on Qwen3.5-9B and GLM-5.3 Flash. We compare it with two prompting floors and a program-induction method, and with the weights of a student, Qwen3.5-9B, trained to imitate a teacher, GLM-5.3 Flash \citep{ross2011dagger}. In some runs GLM-5.3 Flash also writes Qwen3.5-9B's notes, so that notes and weights are taught by the same teacher. On returning tasks, we find that the harness saves less than a floor that simply caches each task's demonstrations, with either model, and only the teacher's harness keeps up with the program-induction method, when explicitly told to store one program per task (Figure~\ref{fig:plane} and Section~\ref{sec:results}). On unseen tasks, notes and programs save nothing measurable, while training on the teacher's moves makes the student considerably better.

We also define collective regimes in which one system serves 2 or 4 users at once, each operating on non-overlapping tasks of their own (Section~\ref{sec:regimes}). This simulates a working environment in which people are responsible for different tasks and the learning is collective over the group. More concurrent users generate more varied training data over the stream's 200 rounds, and our results show that this leads to better generalization to unseen tasks. In contrast, no harness-based method gains on the unseen tasks whatsoever. Our benchmark highlights significant shortcomings in how much off-the-shelf harnesses can ``learn'' in unfamiliar domains, and opens the door to research on training models to drive harnesses in such domains.
